\Needspace{8\baselineskip}
\Article{280}
Nehmen wir an, die binäre Form
\[
att+2btu+cuu\ldots\varphi
\]
\SourceLeaf{324}
mit der Determinante $=D$ werde durch die ternäre Form $f$ mit den
Unbestimmten $x,x',x''$ dargestellt, indem man
\[
x=mt+nu,\qquad x'=m't+n'u,\qquad x''=m''t+n''u
\]
setzt, und der Form $f$ sei die Form $F$ mit den Unbestimmten $X,X',X''$
adjungiert. Dann bestätigt eine Rechnung leicht (wenn man die Koeffizienten
der Formen $f,F$ durch eigene Buchstaben bezeichnet), oder es folgt auch
unmittelbar aus Art.\ 268.\ II., dass die Zahl $D$ durch $F$
dargestellt wird, indem man
\[
X=m'n''-m''n',\qquad X'=m''n-mn'',\qquad X''=mn'-m'n
\]
setzt. Diese Darstellung der Zahl $D$ kann man zweckmäßig der Darstellung
der Form $\varphi$ durch $f$ \textit{adjungiert} nennen. Haben die Werte
von $X,X',X''$ keinen gemeinsamen Teiler, so nennen wir diese Darstellung
von $D$ der Kürze halber \textit{eigentlich}, andernfalls \textit{uneigentlich}.
Dieselben Bezeichnungen geben wir auch der Darstellung der Form $\varphi$
durch $f$, der jene Darstellung von $D$ adjungiert ist. Das Auffinden
aller eigentlichen Darstellungen der Zahl $D$ durch die Form $F$
beruht nun auf folgenden Punkten:

I.\ Jede Darstellung von $D$ durch $F$ lässt sich aus einer
Darstellung irgendeiner Form mit der Determinante $D$ durch die Form
$f$ ableiten, \textit{d.h.}, sie ist einer solchen Darstellung adjungiert.

Sei nämlich irgendeine Darstellung von $D$ durch $F$ durch $X=L$,
$X'=L'$, $X''=L''$ gegeben. Nach dem Lemma des vorigen Artikels wähle man
$m,m',m'',n,n',n''$ so, dass
\[
m'n''-m''n'=L,\qquad m''n-mn''=L',\qquad mn'-m'n=L''
\]
gilt, und die Form $f$ gehe durch die Substitution
\[
x=mt+nu,\qquad x'=m't+n'u,\qquad x''=m''t+n''u
\]
in die binäre Form $\varphi=att+2btu+cuu$ über. Dann sieht man leicht, dass $D$ die
Determinante der Form $\varphi$ ist und dass ihrer Darstellung durch $f$
die vorgelegte Darstellung von $D$ durch $F$ adjungiert ist.

\textit{Beispiel.} Sei $f=xx+x'x'+x''x''$, also $F=-XX-X'X'-X''X''$;
ferner $D=-209$, und ihre Darstellung durch $F$ sei durch $X=1$, $X'=8$, $X''=12$ gegeben.
\SourceLeaf{325}
Daraus ergeben sich für $m,m',m'',n,n',n''$ die Werte $-20,1,1,-12,0,1$, in dieser Reihenfolge,
und $\varphi=402tt+482tu+145uu$.

II.\ Sind $\varphi,\chi$ eigentlich äquivalente binäre Formen, so ist jede
Darstellung von $D$ durch $F$, die einer Darstellung der Form
$\varphi$ durch $f$ adjungiert ist, auch einer Darstellung der Form
$\chi$ durch $f$ adjungiert.

Seien $p,q$ die Unbestimmten der Form $\chi$; die Form $\varphi$ gehe
in $\chi$ durch die eigentliche Substitution $t=\alpha p+\beta q$, $u=\gamma p+\delta q$ über,
und irgendeine Darstellung der Form $\varphi$ durch $f$ sei
\[
x=mt+nu,\qquad x'=m't+n'u,\qquad x''=m''t+n''u\ldots(R)
\]
Dann sieht man ohne Schwierigkeit: Setzt man
\[
\begin{aligned}
\alpha m+\gamma n&=g,&\quad \alpha m'+\gamma n'&=g',&\quad \alpha m''+\gamma n''&=g''\\
\beta m+\delta n&=h,&\quad \beta m'+\delta n'&=h',&\quad \beta m''+\delta n''&=h''
\end{aligned}
\]
so wird die Form $\chi$ durch $f$ dargestellt, indem man
\[
x=gp+hq,\qquad x'=g'p+h'q,\qquad x''=g''p+h''q\ldots(R')
\]
setzt. Durch Rechnung findet man (wegen $\alpha\delta-\beta\gamma=1$)
\[
g'h''-g''h'=m'n''-m''n',\quad
g''h-gh''=m''n-mn'',\quad
gh'-g'h=mn'-m'n
\]
\textit{d.h.}, den Darstellungen $R,R'$ ist dieselbe Darstellung von
$D$ durch $F$ adjungiert.

So ist im vorigen Beispiel die Form $\varphi$ äquivalent zu
$\chi=13pp-10pq+18qq$, in die sie durch die eigentliche Substitution
$t=-3p+q$, $u=5p-2q$ übergeht. Daraus ergibt sich die Darstellung der Form
$\chi$ durch $f$ als $x=4q$, $x'=-3p+q$, $x''=2p-q$, aus der wieder dieselbe
Darstellung der Zahl $-209$ folgt, von der wir ausgegangen waren.

III.\ Können schließlich zwei binäre Formen $\varphi,\chi$ mit der Determinante
$D$ und den Unbestimmten $t,u$; $p,q$ durch $f$ dargestellt
werden, und ist einer Darstellung der einen dieselbe eigentliche Darstellung
von $D$ durch $F$ adjungiert wie einer Darstellung der anderen,
so sind diese Formen notwendig eigentlich äquivalent. Nehmen wir an,
$\varphi$ werde durch $f$ dargestellt, indem man
\[
x=mt+nu,\qquad x'=m't+n'u,\qquad x''=m''t+n''u
\]
\SourceLeaf{326}
setzt, $\chi$ dagegen, indem man
\[
x=gp+hq.\qquad x'=g'p+h'q,\qquad x''=g''p+h''q
\]
setzt, und es gelte
\[
\begin{aligned}
m'n''-m''n'&=g'h''-g''h'=L\\
m''n-mn''&=g''h-gh''=L'\\
mn'-m'n&=gh'-g'h=L''
\end{aligned}
\]
Man wähle ganze Zahlen $l,l',l''$ so, dass $Ll+L'l'+L''l''=1$ gilt, und setze
\[
\begin{aligned}
n'l''-n''l'&=M,&\quad n''l-nl''&=M',&\quad nl'-n'l&=M''\\
l'm''-l''m'&=N,&\quad l''m-lm''&=N',&\quad lm'-l'm&=N''
\end{aligned}
\]
Schließlich setze man
\[
\begin{aligned}
gM+g'M'+g''M''&=\alpha,&\quad hM+h'M'+h''M''&=\beta\\
gN+g'N'+g''N''&=\gamma,&\quad hN+h'N'+h''N''&=\delta
\end{aligned}
\]
Daraus folgt leicht
\[
\begin{aligned}
\alpha m+\gamma n&=g-l(gL+g'L'+g''L'')=g\\
\beta m+\delta n&=h-l(hL+h'L'+h''L'')=h
\end{aligned}
\]
und ebenso
\[
\alpha m'+\gamma n'=g',\quad \beta m'+\delta n'=h',\quad
\alpha m''+\gamma n''=g'',\quad \beta m''+\delta n''=h''
\]
Daraus folgt. $mt+nu$, $m't+n'u$, $m''t+n''u$ gehen durch die Substitution
\[
t=\alpha p+\beta q,\quad u=\gamma p+\delta q\ldots(S)
\]
in $gp+hq$, $g'p+h'q$, $g''p+h''q$ über. Also geht $\varphi$ offenbar durch
die Substitution $S$ in dieselbe Form über wie $f$, wenn man
\[
x=gp+hq,\qquad x'=g'p+h'q,\qquad x''=g''p+h''q
\]
setzt, also in die Form $\chi$, der sie somit äquivalent ist.
Schließlich erhält man durch die nötigen Einsetzungen leicht
\[
\alpha\delta-\beta\gamma=(Ll+L'l'+L''l'')^2=1
\]
Die Substitution $S$ ist also eigentlich, und die Formen $\varphi,\chi$
sind eigentlich äquivalent.
\SourceLeaf{327}

Aus diesen Bemerkungen folgen die folgenden Regeln, um alle eigentlichen
Darstellungen von $D$ durch $F$ zu finden: Man bestimme alle
Klassen binärer Formen mit der Determinante $D$ und wähle aus jeder
eine beliebige Form. Man suche alle eigentlichen Darstellungen jeder dieser
Formen durch $f$ (wobei man diejenigen verwirft, die nicht durch
$f$ dargestellt werden können) und leite aus jeder dieser Darstellungen
die Darstellungen der Zahl $D$ durch $F$ ab. Aus I und II ist
offenbar, dass man so alle möglichen eigentlichen Darstellungen erhält,
die Lösung also vollständig ist; aus III, dass Transformationen von Formen
aus verschiedenen Klassen sicher verschiedene Darstellungen liefern.

\Needspace{8\baselineskip}
\Article{281}
Das Auffinden der \textit{uneigentlichen} Darstellungen einer gegebenen Zahl
$D$ durch die Form $F$ lässt sich leicht auf den vorigen Fall
zurückführen. Ist nämlich $D$ durch kein Quadrat (außer $1$)
teilbar, so gibt es offenbar überhaupt keine solchen Darstellungen.
Andernfalls, wenn $D$ durch die Quadrate $\lambda\lambda,\mu\mu,\nu\nu$ usw.\ teilbar ist,
findet man alle uneigentlichen Darstellungen von $D$ durch $F$,
indem man alle eigentlichen Darstellungen der Zahlen $\frac{D}{\lambda\lambda}$,
$\frac{D}{\mu\mu}$, $\frac{D}{\nu\nu}$ usw.\ durch dieselbe Form bestimmt und die Werte der
Unbestimmten jeweils mit $\lambda,\mu,\nu$ usw.\ multipliziert.

So hängt das Auffinden aller Darstellungen einer gegebenen Zahl durch eine
gegebene ternäre Form, \textit{die einer ternären Form adjungiert ist},
vom zweiten Problem ab. Auf diesen Fall, der auf den ersten Blick weniger
weit zu reichen scheint, lassen sich die übrigen folgendermaßen zurückführen.
Sei $D$ die Zahl, die durch die Form
$\left(\begin{smallmatrix}g,g',g''\\h,h',h''\end{smallmatrix}\right)$
mit der Determinante $\Delta$ dargestellt werden soll; dieser sei die Form
$\left(\begin{smallmatrix}G,G',G''\\H,H',H''\end{smallmatrix}\right)=f$
adjungiert. Dieser wiederum ist die Form
$\left(\begin{smallmatrix}\Delta g,\Delta g',\Delta g''\\\Delta h,\Delta h',\Delta h''\end{smallmatrix}\right)=F$
adjungiert, und offenbar sind die Darstellungen der Zahl $\Delta D$ durch
$F$ (deren Bestimmung vom Vorangegangenen abhängt) genau dieselben
wie die Darstellungen der Zahl $D$ durch die vorgelegte Form.
--- Haben ferner alle Koeffizienten der Form $f$ einen gemeinsamen
Teiler $\mu$, so sind offenbar alle Koeffizienten der Form $F$
durch $\mu\mu$ teilbar. Also muss auch $\Delta D$ durch $\mu\mu$ teilbar
sein (andernfalls gäbe es keine Darstellungen); und die Darstellungen
der Zahl $D$ durch die vorgelegte Form stimmen mit den Darstellungen
der Zahl $\frac{\Delta D}{\mu\mu}$ durch die Form überein, die aus $F$ entsteht,
indem man jeden Koeffizienten durch $\mu\mu$ teilt. Diese Form ist derjenigen
adjungiert, die aus $f$ entsteht, indem man jeden Koeffizienten
durch $\mu$ teilt.

Schließlich bemerken wir, dass diese Lösung des ersten Problems in dem
einzigen Fall, wo
\SourceLeaf{328}
$D=0$ ist, nicht anwendbar ist; denn hier verteilen sich nicht alle
binären Formen mit der Determinante $D$ auf eine endliche Anzahl von
Klassen. Diesen Fall werden wir unten nach anderen Grundsätzen lösen.

